%
% claude-code.tex
%
% Z Specification for Claude Code State Machine
% Formal model of the session lifecycle, agentic loop, tool execution
% pipeline, subagent coordination, and hook system.
%
% Type-check: fuzz -t claude-code.tex
% Animate:    probcli claude-code.tex -animate 20
%

\documentclass[a4paper,10pt,fleqn]{article}
\usepackage[margin=1in]{geometry}
\usepackage{fuzz}
\usepackage[colorlinks=true,linkcolor=blue,citecolor=blue,urlcolor=blue]{hyperref}

\begin{document}

\title{Claude Code: A Z Specification}
\author{Formal Model of the Agentic Session Lifecycle}
\date{March 2026}
\hypersetup{
  pdftitle={Claude Code: A Z Specification},
  pdfauthor={Punt Labs},
  pdfsubject={Z Specification},
  pdfcreator={fuzz/probcli}
}
\maketitle

\tableofcontents
\newpage

% ============================================================
\section{Introduction}
% ============================================================

This specification models the state machine of Claude Code, Anthropic's
agentic coding assistant. The model captures three concurrent state
dimensions:

\begin{enumerate}
  \item \textbf{Session phase} --- the top-level lifecycle from startup
    through the agentic loop to session end.
  \item \textbf{Tool execution phase} --- the pipeline a single tool call
    traverses: pre-hook, permission check, execution, post-hook.
  \item \textbf{Subagent phase} --- spawning, running, and stopping of
    subordinate agents.
\end{enumerate}

Hooks are modeled as \emph{guard conditions} on transitions. A blocking
hook (one that can prevent a transition) appears as a precondition input
on the operation schema. A non-blocking hook (observation only) is noted
in the narrative but does not appear in the formal model because it
cannot alter the state machine's reachable states.

The payoff: once you know the precise states and transitions, you can
systematically attach workflow requirements (commit discipline, lint
gates, permission policies) to the correct hook points with confidence
that no path through the machine escapes your controls.

% ============================================================
\section{Basic Types}
% ============================================================

\begin{zed}
  [TOOLNAME, TOOLINPUT, TOOLOUTPUT, PROMPT, RESPONSE, AGENTID, SESSIONID, CONTEXTCHUNK]
\end{zed}

\texttt{TOOLNAME} identifies a tool (e.g., \texttt{Bash}, \texttt{Edit},
\texttt{mcp\_\_github\_\_create\_pull\_request}). \texttt{TOOLINPUT} and
\texttt{TOOLOUTPUT} are the opaque input and output payloads of a tool
call. \texttt{PROMPT} is a user-submitted prompt. \texttt{RESPONSE} is
Claude's generated text. \texttt{AGENTID} uniquely identifies a spawned
subagent. \texttt{SESSIONID} identifies a session. \texttt{CONTEXTCHUNK}
is a unit of conversation context (message, tool result, etc.).

% ============================================================
\section{Free Types}
% ============================================================

\begin{zed}
  SessionPhase ::= spInactive | spIdle | spProcessing | spResponding | spEnded
\end{zed}

The session lifecycle. \texttt{spInactive} is before \texttt{SessionStart}.
\texttt{spIdle} is waiting for user input. \texttt{spProcessing} is Claude
reasoning and dispatching tools. \texttt{spResponding} is streaming the
final response to the user. \texttt{spEnded} is terminal. Startup and
shutdown hooks are modeled as atomic transitions (not intermediate states)
because no operation is enabled during hook execution.

\begin{zed}
  ToolPhase ::= \\
  \quad~ tpNone | tpPreHook | tpPermission | tpExecuting | tpPostHook
\end{zed}

The tool execution pipeline. \texttt{tpNone} means no tool call is
active. \texttt{tpPreHook} is running \texttt{PreToolUse} hooks.
\texttt{tpPermission} is awaiting user permission (if required).
\texttt{tpExecuting} is the tool running. \texttt{tpPostHook} is
running \texttt{PostToolUse} or \texttt{PostToolUseFailure} hooks.

\begin{zed}
  PermissionMode ::= \\
  \quad~ pmDefault | pmPlan | pmAcceptEdits | pmDontAsk | pmBypass
\end{zed}

The five permission modes. \texttt{pmDefault} prompts for dangerous
operations. \texttt{pmPlan} is read-only. \texttt{pmAcceptEdits}
auto-approves edits. \texttt{pmDontAsk} auto-approves everything.
\texttt{pmBypass} disables all permission checks.

\begin{zed}
  PermissionDecision ::= pdAllow | pdDeny | pdAsk
\end{zed}

A \texttt{PreToolUse} hook can allow, deny, or escalate to the user.

\begin{zed}
  HookDecision ::= hdAllow | hdBlock
\end{zed}

For blocking hooks on \texttt{Stop}, \texttt{SubagentStop},
\texttt{UserPromptSubmit}, \texttt{TaskCompleted}, \texttt{TeammateIdle},
and \texttt{ConfigChange}: the hook either allows the transition or
blocks it.

\begin{zed}
  SessionStartSource ::= ssStartup | ssResume | ssClear | ssCompact
\end{zed}

How the session was initiated.

\begin{zed}
  SessionEndReason ::= \\
  \quad~ seClear | seLogout | sePromptExit | seBypassDisabled | seOther
\end{zed}

Why the session ended.

\begin{zed}
  ToolResult ::= trSuccess | trFailure
\end{zed}

Whether a tool execution succeeded or failed, determining which
post-hook fires (\texttt{PostToolUse} vs.\ \texttt{PostToolUseFailure}).

\begin{zed}
  ZBOOL ::= ztrue | zfalse
\end{zed}

% ============================================================
\section{Global Constants}
% ============================================================

\begin{axdef}
  maxTools : \nat \\
  maxSubagents : \nat \\
  maxContext : \nat
  \where
  maxTools = 20 \\
  maxSubagents = 10 \\
  maxContext = 200
\end{axdef}

Bounds for ProB animation. \texttt{maxTools} caps tools dispatched per
turn. \texttt{maxSubagents} caps concurrent subagents.
\texttt{maxContext} caps context window chunks.

% ============================================================
\section{State}
% ============================================================

\begin{schema}{State}
  sessionId : \finset SESSIONID \\
  sessionPhase : SessionPhase \\
  permMode : PermissionMode \\
  %
  toolPhase : ToolPhase \\
  currentTool : \finset TOOLNAME \\
  toolsThisTurn : \nat \\
  %
  activeSubagents : \finset AGENTID \\
  %
  context : \seq CONTEXTCHUNK \\
  pendingPrompt : \finset PROMPT \\
  pendingResponse : \finset RESPONSE
  \where
  %
  % Session identity: at most one active session
  %
  \# sessionId \leq 1 \\
  sessionPhase = spInactive \implies sessionId = \emptyset \\
  sessionPhase \neq spInactive \implies \# sessionId = 1 \\
  %
  % Tool phase is only active during processing
  %
  sessionPhase \neq spProcessing \implies toolPhase = tpNone \\
  %
  % Current tool tracks what is executing
  %
  \# currentTool \leq 1 \\
  toolPhase = tpNone \implies currentTool = \emptyset \\
  toolPhase \neq tpNone \implies \# currentTool = 1 \\
  %
  % Tool count bounded for animation
  %
  toolsThisTurn \leq maxTools \\
  %
  % Subagents only exist during active session
  %
  sessionPhase = spInactive \implies activeSubagents = \emptyset \\
  sessionPhase = spEnded \implies activeSubagents = \emptyset \\
  \# activeSubagents \leq maxSubagents \\
  %
  % Context bounded for animation
  %
  \# context \leq maxContext \\
  %
  % Pending prompt: at most one, only when processing
  %
  \# pendingPrompt \leq 1 \\
  sessionPhase = spIdle \implies pendingPrompt = \emptyset \\
  sessionPhase = spInactive \implies pendingPrompt = \emptyset \\
  sessionPhase = spEnded \implies pendingPrompt = \emptyset \\
  %
  % Pending response: at most one, only when responding
  %
  \# pendingResponse \leq 1 \\
  sessionPhase \neq spResponding \implies pendingResponse = \emptyset
\end{schema}

The state is \emph{flat}: session, tool, and subagent dimensions are
composed in a single schema so that cross-cutting invariants (e.g.,
``tool phase is \texttt{tpNone} unless session is processing'') are
enforced directly.

% ============================================================
\section{Initialization}
% ============================================================

\begin{schema}{Init}
  State'
  \where
  sessionId' = \emptyset \\
  sessionPhase' = spInactive \\
  permMode' = pmDefault \\
  toolPhase' = tpNone \\
  currentTool' = \emptyset \\
  toolsThisTurn' = 0 \\
  activeSubagents' = \emptyset \\
  context' = \langle \rangle \\
  pendingPrompt' = \emptyset \\
  pendingResponse' = \emptyset
\end{schema}

% ============================================================
\section{Session Lifecycle Operations}
% ============================================================

% ------------------------------------------------------------
\subsection{StartSession}
% ------------------------------------------------------------

Transitions from \texttt{spInactive} to \texttt{spIdle}. The
\texttt{SessionStart} hook fires during this transition (modeled
as the \texttt{spStarting} intermediate state). \texttt{SessionStart}
cannot block --- it is observational.

\begin{schema}{StartSession}
  \Delta State \\
  sid? : SESSIONID \\
  source? : SessionStartSource \\
  mode? : PermissionMode
  \where
  sessionPhase = spInactive \\
  %
  sessionPhase' = spIdle \\
  sessionId' = \{ sid? \} \\
  permMode' = mode? \\
  toolPhase' = tpNone \\
  currentTool' = \emptyset \\
  toolsThisTurn' = 0 \\
  activeSubagents' = \emptyset \\
  context' = \langle \rangle \\
  pendingPrompt' = \emptyset \\
  pendingResponse' = \emptyset
\end{schema}

Hook event: \texttt{SessionStart} fires with \texttt{source?} as the
matcher value. The hook can inject \texttt{additionalContext} into the
context window and persist environment variables via
\texttt{CLAUDE\_ENV\_FILE}. It cannot block.

% ------------------------------------------------------------
\subsection{SubmitPrompt}
% ------------------------------------------------------------

The user submits a prompt. \texttt{UserPromptSubmit} fires and can
block (exit 2 erases the prompt).

\begin{schema}{SubmitPrompt}
  \Delta State \\
  prompt? : PROMPT \\
  hookResult? : HookDecision \\
  contextChunk? : CONTEXTCHUNK
  \where
  sessionPhase = spIdle \\
  hookResult? = hdAllow \\
  %
  sessionPhase' = spProcessing \\
  pendingPrompt' = \{ prompt? \} \\
  context' = (1 \upto maxContext) \dres (context \cat \langle contextChunk? \rangle) \\
  toolsThisTurn' = 0 \\
  toolPhase' = tpNone \\
  currentTool' = \emptyset \\
  %
  % Frame: unchanged
  %
  sessionId' = sessionId \\
  permMode' = permMode \\
  activeSubagents' = activeSubagents \\
  pendingResponse' = \emptyset
\end{schema}

If \texttt{hookResult?} were \texttt{hdBlock}, the precondition fails
and the operation is refused --- the prompt is erased and the session
stays in \texttt{spIdle}.

\begin{schema}{RejectPrompt}
  \Xi State \\
  prompt? : PROMPT \\
  hookResult? : HookDecision
  \where
  sessionPhase = spIdle \\
  hookResult? = hdBlock
\end{schema}

% ------------------------------------------------------------
\subsection{BeginToolCall}
% ------------------------------------------------------------

Claude decides to invoke a tool. Enters the \texttt{PreToolUse} hook
phase.

\begin{schema}{BeginToolCall}
  \Delta State \\
  tool? : TOOLNAME
  \where
  sessionPhase = spProcessing \\
  toolPhase = tpNone \\
  toolsThisTurn < maxTools \\
  %
  toolPhase' = tpPreHook \\
  currentTool' = \{ tool? \} \\
  %
  % Frame: unchanged
  %
  sessionPhase' = sessionPhase \\
  sessionId' = sessionId \\
  permMode' = permMode \\
  toolsThisTurn' = toolsThisTurn \\
  activeSubagents' = activeSubagents \\
  context' = context \\
  pendingPrompt' = pendingPrompt \\
  pendingResponse' = pendingResponse
\end{schema}

% ------------------------------------------------------------
\subsection{PreToolHookAllow}
% ------------------------------------------------------------

The \texttt{PreToolUse} hook allows the tool call. If the permission
mode requires a check, we move to \texttt{tpPermission}. Otherwise,
we skip directly to \texttt{tpExecuting}.

\begin{schema}{PreToolHookAllow}
  \Delta State \\
  hookDecision? : PermissionDecision \\
  needsPermission? : ZBOOL
  \where
  sessionPhase = spProcessing \\
  toolPhase = tpPreHook \\
  hookDecision? = pdAllow \\
  %
  toolPhase' = \IF needsPermission? = ztrue \THEN tpPermission \ELSE tpExecuting \\
  %
  % Frame: unchanged
  %
  sessionPhase' = sessionPhase \\
  sessionId' = sessionId \\
  permMode' = permMode \\
  currentTool' = currentTool \\
  toolsThisTurn' = toolsThisTurn \\
  activeSubagents' = activeSubagents \\
  context' = context \\
  pendingPrompt' = pendingPrompt \\
  pendingResponse' = pendingResponse
\end{schema}

% ------------------------------------------------------------
\subsection{PreToolHookDeny}
% ------------------------------------------------------------

The \texttt{PreToolUse} hook denies the tool call. The tool phase
resets to \texttt{tpNone} and Claude receives the denial reason.

\begin{schema}{PreToolHookDeny}
  \Delta State \\
  hookDecision? : PermissionDecision \\
  denialContext? : CONTEXTCHUNK
  \where
  sessionPhase = spProcessing \\
  toolPhase = tpPreHook \\
  hookDecision? = pdDeny \\
  %
  toolPhase' = tpNone \\
  currentTool' = \emptyset \\
  context' = (1 \upto maxContext) \dres (context \cat \langle denialContext? \rangle) \\
  %
  % Frame: unchanged
  %
  sessionPhase' = sessionPhase \\
  sessionId' = sessionId \\
  permMode' = permMode \\
  toolsThisTurn' = toolsThisTurn \\
  activeSubagents' = activeSubagents \\
  pendingPrompt' = pendingPrompt \\
  pendingResponse' = pendingResponse
\end{schema}

% ------------------------------------------------------------
\subsection{PreToolHookEscalate}
% ------------------------------------------------------------

The \texttt{PreToolUse} hook returns \texttt{pdAsk}, escalating to
the permission dialog regardless of permission mode.

\begin{schema}{PreToolHookEscalate}
  \Delta State \\
  hookDecision? : PermissionDecision
  \where
  sessionPhase = spProcessing \\
  toolPhase = tpPreHook \\
  hookDecision? = pdAsk \\
  %
  toolPhase' = tpPermission \\
  %
  % Frame: unchanged
  %
  sessionPhase' = sessionPhase \\
  sessionId' = sessionId \\
  permMode' = permMode \\
  currentTool' = currentTool \\
  toolsThisTurn' = toolsThisTurn \\
  activeSubagents' = activeSubagents \\
  context' = context \\
  pendingPrompt' = pendingPrompt \\
  pendingResponse' = pendingResponse
\end{schema}

% ------------------------------------------------------------
\subsection{GrantPermission}
% ------------------------------------------------------------

The user (or \texttt{PermissionRequest} hook) grants permission.
The \texttt{PermissionRequest} hook fires during this phase and
can allow or deny.

\begin{schema}{GrantPermission}
  \Delta State
  \where
  sessionPhase = spProcessing \\
  toolPhase = tpPermission \\
  %
  toolPhase' = tpExecuting \\
  %
  % Frame: unchanged
  %
  sessionPhase' = sessionPhase \\
  sessionId' = sessionId \\
  permMode' = permMode \\
  currentTool' = currentTool \\
  toolsThisTurn' = toolsThisTurn \\
  activeSubagents' = activeSubagents \\
  context' = context \\
  pendingPrompt' = pendingPrompt \\
  pendingResponse' = pendingResponse
\end{schema}

% ------------------------------------------------------------
\subsection{DenyPermission}
% ------------------------------------------------------------

The user (or \texttt{PermissionRequest} hook) denies permission.

\begin{schema}{DenyPermission}
  \Delta State \\
  denialContext? : CONTEXTCHUNK
  \where
  sessionPhase = spProcessing \\
  toolPhase = tpPermission \\
  %
  toolPhase' = tpNone \\
  currentTool' = \emptyset \\
  context' = (1 \upto maxContext) \dres (context \cat \langle denialContext? \rangle) \\
  %
  % Frame: unchanged
  %
  sessionPhase' = sessionPhase \\
  sessionId' = sessionId \\
  permMode' = permMode \\
  toolsThisTurn' = toolsThisTurn \\
  activeSubagents' = activeSubagents \\
  pendingPrompt' = pendingPrompt \\
  pendingResponse' = pendingResponse
\end{schema}

% ------------------------------------------------------------
\subsection{CompleteToolSuccess}
% ------------------------------------------------------------

The tool executes successfully. \texttt{PostToolUse} fires (non-blocking).

\begin{schema}{CompleteToolSuccess}
  \Delta State \\
  output? : CONTEXTCHUNK
  \where
  sessionPhase = spProcessing \\
  toolPhase = tpExecuting \\
  %
  toolPhase' = tpNone \\
  currentTool' = \emptyset \\
  toolsThisTurn' = toolsThisTurn + 1 \\
  context' = (1 \upto maxContext) \dres (context \cat \langle output? \rangle) \\
  %
  % Frame: unchanged
  %
  sessionPhase' = sessionPhase \\
  sessionId' = sessionId \\
  permMode' = permMode \\
  activeSubagents' = activeSubagents \\
  pendingPrompt' = pendingPrompt \\
  pendingResponse' = pendingResponse
\end{schema}

% ------------------------------------------------------------
\subsection{CompleteToolFailure}
% ------------------------------------------------------------

The tool fails. \texttt{PostToolUseFailure} fires (non-blocking).

\begin{schema}{CompleteToolFailure}
  \Delta State \\
  errorContext? : CONTEXTCHUNK
  \where
  sessionPhase = spProcessing \\
  toolPhase = tpExecuting \\
  %
  toolPhase' = tpNone \\
  currentTool' = \emptyset \\
  toolsThisTurn' = toolsThisTurn + 1 \\
  context' = (1 \upto maxContext) \dres (context \cat \langle errorContext? \rangle) \\
  %
  % Frame: unchanged
  %
  sessionPhase' = sessionPhase \\
  sessionId' = sessionId \\
  permMode' = permMode \\
  activeSubagents' = activeSubagents \\
  pendingPrompt' = pendingPrompt \\
  pendingResponse' = pendingResponse
\end{schema}

% ============================================================
\section{Subagent Operations}
% ============================================================

% ------------------------------------------------------------
\subsection{SpawnSubagent}
% ------------------------------------------------------------

Claude spawns a subagent. \texttt{SubagentStart} fires (non-blocking).

\begin{schema}{SpawnSubagent}
  \Delta State \\
  agentId? : AGENTID
  \where
  sessionPhase = spProcessing \\
  agentId? \notin activeSubagents \\
  \# activeSubagents < maxSubagents \\
  %
  activeSubagents' = activeSubagents \cup \{ agentId? \} \\
  %
  % Frame: unchanged
  %
  sessionPhase' = sessionPhase \\
  sessionId' = sessionId \\
  permMode' = permMode \\
  toolPhase' = toolPhase \\
  currentTool' = currentTool \\
  toolsThisTurn' = toolsThisTurn \\
  context' = context \\
  pendingPrompt' = pendingPrompt \\
  pendingResponse' = pendingResponse
\end{schema}

% ------------------------------------------------------------
\subsection{StopSubagent}
% ------------------------------------------------------------

A subagent finishes. \texttt{SubagentStop} fires and can block (exit 2
prevents the subagent from stopping, forcing it to continue).

\begin{schema}{StopSubagent}
  \Delta State \\
  agentId? : AGENTID \\
  hookResult? : HookDecision \\
  resultContext? : CONTEXTCHUNK
  \where
  sessionPhase = spProcessing \\
  agentId? \in activeSubagents \\
  hookResult? = hdAllow \\
  %
  activeSubagents' = activeSubagents \setminus \{ agentId? \} \\
  context' = (1 \upto maxContext) \dres (context \cat \langle resultContext? \rangle) \\
  %
  % Frame: unchanged
  %
  sessionPhase' = sessionPhase \\
  sessionId' = sessionId \\
  permMode' = permMode \\
  toolPhase' = toolPhase \\
  currentTool' = currentTool \\
  toolsThisTurn' = toolsThisTurn \\
  pendingPrompt' = pendingPrompt \\
  pendingResponse' = pendingResponse
\end{schema}

% ============================================================
\section{Response and Stop Operations}
% ============================================================

% ------------------------------------------------------------
\subsection{BeginResponse}
% ------------------------------------------------------------

Claude has finished tool calls and begins streaming a response.

\begin{schema}{BeginResponse}
  \Delta State \\
  response? : RESPONSE
  \where
  sessionPhase = spProcessing \\
  toolPhase = tpNone \\
  activeSubagents = \emptyset \\
  %
  sessionPhase' = spResponding \\
  pendingResponse' = \{ response? \} \\
  pendingPrompt' = \emptyset \\
  %
  % Frame: unchanged
  %
  sessionId' = sessionId \\
  permMode' = permMode \\
  toolPhase' = tpNone \\
  currentTool' = \emptyset \\
  toolsThisTurn' = toolsThisTurn \\
  activeSubagents' = activeSubagents \\
  context' = context
\end{schema}

% ------------------------------------------------------------
\subsection{FinishResponse}
% ------------------------------------------------------------

Claude finishes responding. The \texttt{Stop} hook fires and can block
(exit 2 prevents stopping, forcing Claude to continue processing).

\begin{schema}{FinishResponse}
  \Delta State \\
  hookResult? : HookDecision \\
  responseContext? : CONTEXTCHUNK
  \where
  sessionPhase = spResponding \\
  hookResult? = hdAllow \\
  %
  sessionPhase' = spIdle \\
  pendingResponse' = \emptyset \\
  context' = (1 \upto maxContext) \dres (context \cat \langle responseContext? \rangle) \\
  toolsThisTurn' = 0 \\
  %
  % Frame: unchanged
  %
  sessionId' = sessionId \\
  permMode' = permMode \\
  toolPhase' = tpNone \\
  currentTool' = \emptyset \\
  activeSubagents' = activeSubagents \\
  pendingPrompt' = \emptyset
\end{schema}

% ------------------------------------------------------------
\subsection{ForceContiue}
% ------------------------------------------------------------

The \texttt{Stop} hook blocks the stop. Claude returns to processing.

\begin{schema}{ForceContinue}
  \Delta State \\
  hookResult? : HookDecision
  \where
  sessionPhase = spResponding \\
  hookResult? = hdBlock \\
  %
  sessionPhase' = spProcessing \\
  pendingResponse' = \emptyset \\
  %
  % Frame: unchanged
  %
  sessionId' = sessionId \\
  permMode' = permMode \\
  toolPhase' = tpNone \\
  currentTool' = \emptyset \\
  toolsThisTurn' = toolsThisTurn \\
  activeSubagents' = activeSubagents \\
  context' = context \\
  pendingPrompt' = pendingPrompt
\end{schema}

% ============================================================
\section{Context Management}
% ============================================================

% ------------------------------------------------------------
\subsection{CompactContext}
% ------------------------------------------------------------

Context compaction occurs when the context window approaches its limit.
\texttt{PreCompact} fires (non-blocking). The context is compressed.
After compaction, \texttt{SessionStart} fires with \texttt{ssCompact}
as the source.

\begin{schema}{CompactContext}
  \Delta State \\
  newLen? : \nat
  \where
  sessionPhase = spIdle \\
  \# context > 0 \\
  newLen? < \# context \\
  newLen? \leq maxContext \\
  %
  context' = (1 \upto newLen?) \dres context \\
  %
  % Frame: unchanged
  %
  sessionPhase' = sessionPhase \\
  sessionId' = sessionId \\
  permMode' = permMode \\
  toolPhase' = toolPhase \\
  currentTool' = currentTool \\
  toolsThisTurn' = toolsThisTurn \\
  activeSubagents' = activeSubagents \\
  pendingPrompt' = pendingPrompt \\
  pendingResponse' = pendingResponse
\end{schema}

% ============================================================
\section{Session End}
% ============================================================

% ------------------------------------------------------------
\subsection{EndSession}
% ------------------------------------------------------------

The session terminates. \texttt{SessionEnd} fires (non-blocking).
All subagents must have stopped.

\begin{schema}{EndSession}
  \Delta State \\
  reason? : SessionEndReason
  \where
  sessionPhase = spIdle \\
  activeSubagents = \emptyset \\
  %
  sessionPhase' = spEnded \\
  sessionId' = \emptyset \\
  toolPhase' = tpNone \\
  currentTool' = \emptyset \\
  toolsThisTurn' = 0 \\
  activeSubagents' = \emptyset \\
  context' = \langle \rangle \\
  pendingPrompt' = \emptyset \\
  pendingResponse' = \emptyset \\
  permMode' = permMode
\end{schema}

% ============================================================
\section{Hook Event Summary}
% ============================================================

The following table maps each Claude Code hook event to its
corresponding operation(s) in this specification and whether the hook
can block the transition.

\begin{tabular}{llll}
\hline
\textbf{Hook Event} & \textbf{Operation(s)} & \textbf{Blocks?} & \textbf{Modeled As} \\
\hline
SessionStart & StartSession & No & Transition guard \\
UserPromptSubmit & SubmitPrompt / RejectPrompt & Yes & hookResult? \\
PreToolUse & PreToolHookAllow / Deny / Escalate & Yes & hookDecision? \\
PermissionRequest & GrantPermission / DenyPermission & Yes & Separate ops \\
PostToolUse & CompleteToolSuccess & No & Side effect \\
PostToolUseFailure & CompleteToolFailure & No & Side effect \\
SubagentStart & SpawnSubagent & No & Side effect \\
SubagentStop & StopSubagent & Yes & hookResult? \\
Stop & FinishResponse / ForceContinue & Yes & hookResult? \\
PreCompact & CompactContext & No & Side effect \\
SessionEnd & EndSession & No & Side effect \\
Notification & (not modeled) & No & External \\
TeammateIdle & (not modeled) & Yes & Agent Teams \\
TaskCompleted & (not modeled) & Yes & Task tracking \\
ConfigChange & (not modeled) & Yes & Config system \\
WorktreeCreate & (not modeled) & Yes & Git worktree \\
WorktreeRemove & (not modeled) & No & Git worktree \\
InstructionsLoaded & (not modeled) & No & Config system \\
\hline
\end{tabular}

Events marked ``not modeled'' are orthogonal to the core state machine.
They can be added as refinements if workflow requirements target them.

% ============================================================
\section{System Invariants}
% ============================================================

The specification enforces these properties across all reachable states:

\begin{enumerate}
  \item \textbf{Tool execution requires processing}: tool calls can only
    occur when the session is in \texttt{spProcessing}. This prevents
    tool dispatch during idle, responding, or ended states.

  \item \textbf{Active tool requires tool phase}: a tool name is tracked
    if and only if the tool phase is not \texttt{tpNone}. No ``orphan''
    tool references.

  \item \textbf{Subagents require active session}: subagents cannot exist
    in \texttt{spInactive} or \texttt{spEnded}.

  \item \textbf{Response requires streaming}: a pending response only
    exists during \texttt{spResponding}.

  \item \textbf{Prompt cleared on response}: the pending prompt is
    cleared when Claude begins responding.

  \item \textbf{Session identity}: exactly one session ID exists when
    the session is active; none when inactive or ended.

  \item \textbf{Clean termination}: \texttt{EndSession} requires all
    subagents to have stopped.

  \item \textbf{Blocking hooks are preconditions}: a blocking hook
    prevents a transition by failing the precondition. The operation
    simply does not occur --- no partial state update is possible.
\end{enumerate}

% ============================================================
\section{Workflow Attachment Points}
% ============================================================

This section identifies where workflow standards attach to the state
machine. Each attachment point names the hook event, the operation it
guards, and example workflow requirements.

\subsection{Pre-commit Discipline}

Attach to \texttt{Stop} (FinishResponse). A hook inspects
\texttt{git status} and blocks the stop if there are unstaged changes
that should be committed. This forces Claude to run the session close
protocol before declaring ``done.''

\subsection{Lint and Format Gates}

Attach to \texttt{PostToolUse} with matcher \texttt{Edit|Write}. After
every file modification, run linters and formatters. Non-blocking: inject
lint errors as context so Claude self-corrects.

\subsection{Permission Policies}

Attach to \texttt{PreToolUse} with tool-specific matchers. Deny
destructive commands (\texttt{rm -rf}, \texttt{git push --force}),
enforce read-only in plan mode, require confirmation for external API
calls.

\subsection{Subagent Governance}

Attach to \texttt{SubagentStop}. Validate that subagent output meets
quality criteria before accepting. Block the stop to force the subagent
to continue if output is incomplete.

\subsection{Prompt Validation}

Attach to \texttt{UserPromptSubmit}. Validate that user prompts include
required context (e.g., bead references for tracked work). Reject prompts
that would start untracked work.

\subsection{Session Bookkeeping}

Attach to \texttt{SessionStart} and \texttt{SessionEnd}. On start: load
project context, sync beads, set environment. On end: sync beads, verify
clean git state, push if needed.

\end{document}
